\input{estilo}
% La opción 3.1 de la pizarra: los nodos sueltos; el receive (rcv) y la
% escritura llegan al nodo 1, el primary, que copia a 2 y 3 (en verde). El
% primary está tachado y de él baja "visibility timeout".
\begin{document}
\begin{figura}\begin{tikzpicture}[n/.style={circle, draw=tinta, fill=relleno, minimum size=8mm, inner sep=0pt}]
  \node[n] (p) at (0,0) {1};
  \node[n] at (0.1,1.6) {};
  \node[n] (s2) at (2.1,2.2) {2};
  \node[n] at (3.4,1.0) {};
  \node[n] (s3) at (3.2,-1.1) {3};
  \node[n] at (1.2,-1.5) {};
  \draw[msg] (-3.6,-0.9) -- (-0.5,-0.2);
  \node at (-2.8,1.5) {rcv};
  \draw[msg] (-1.4,1.1) to[out=0, in=130] (-0.35,0.35);
  \draw[msg, draw=segundo] (p) -- (s2);
  \draw[msg, draw=segundo] (p) to[bend right=12] (s3);
  \tachar{p}
  \draw[msg] (0,-0.5) -- (0,-1.1);
  \node[align=center, anchor=north] at (0,-1.15) {visibility\\timeout};
\end{tikzpicture}\end{figura}
\end{document}
